

\definecolor{strategyarrowone}{HTML}{A98EDA}
\definecolor{strategyarrowtwo}{HTML}{E58C9A}
\definecolor{strategyarrowthree}{HTML}{D1A647}
\definecolor{parentlink}{HTML}{B7B2C2}
\definecolor{strategyink}{HTML}{657BBB}
\definecolor{strategyfill}{HTML}{F1F3FB}
\definecolor{figureink}{HTML}{253049}
\definecolor{codegreen}{HTML}{309F86}
\definecolor{codefill}{HTML}{EEF8F4}

\begin{tikzpicture}[
  font=\sffamily, text=figureink,
  line cap=round, line join=round,
  sf arrow/.style={-{Stealth[length=2.15mm,width=1.55mm]},
    line width=0.95pt,shorten >=1.3pt},
  sf skip/.style={sf arrow,
    preaction={draw=white,line width=3pt,arrows=-,shorten <=1.7pt}},
  sf strategy/.style={rectangle,rounded corners=4pt,
    draw=strategyink,fill=strategyfill,line width=1.05pt,
    minimum width=1.65cm,minimum height=0.70cm,
    inner sep=0pt,font=\large},
  sf parent/.style={rectangle,rounded corners=5pt,
    draw=parentlink!65,fill=parentlink!7,line width=0.7pt,
    minimum width=2.05cm,minimum height=0.90cm,
    align=center,inner sep=5pt},
  sf code/.style={circle,draw=codegreen,fill=codefill,
    line width=0.95pt,minimum size=0.43cm,inner sep=0pt}
]
  % One consistent grid; edit the x coordinates to adjust overall spacing.
  \node[sf parent] (prompt) at (0,0)
    {\textbf{Parent}\\[-1pt]{\small context $x_p$}};
  \foreach \j/\x in {1/2.70,2/5.35,3/8.00,4/10.65}{
    \node[sf strategy] (s\j) at (\x,0) {$s_{p,\j}$};
  }

  % Quiet parent bus. Distinct ports keep arrowheads separated.
  \coordinate (bus) at (0,2.52);
  \coordinate (p2) at ($(s2.north)+(-0.52,0)$);
  \coordinate (p3) at ($(s3.north)+(-0.52,0)$);
  \coordinate (p4) at ($(s4.north)+(0.56,0)$);
  \draw[sf arrow,draw=parentlink] (prompt.east) -- (s1.west);
  \draw[draw=parentlink,line width=0.75pt,rounded corners=4pt]
    (prompt.north) -- (bus) -- (bus -| p4);
  \foreach \pin in {p2,p3,p4}{
    \draw[sf arrow,draw=parentlink,line width=0.75pt]
      (bus -| \pin) -- (\pin);
  }

  % Adjacent strategy dependencies; colors encode the source strategy.
  \draw[sf arrow,draw=strategyarrowone] (s1.east) -- (s2.west);
  \draw[sf arrow,draw=strategyarrowtwo] (s2.east) -- (s3.west);
  \draw[sf arrow,draw=strategyarrowthree] (s3.east) -- (s4.west);

  % Long-range dependencies. A narrow white under-stroke separates crossings;
  % replace 'white' in sf skip if your slide background has another color.
  \draw[sf skip,draw=strategyarrowone]
    ($(s1.north)+(-0.28,0)$)
    .. controls +(0,2.18) and +(0,2.18) ..
    ($(s4.north)+(0.22,0)$);
  \draw[sf skip,draw=strategyarrowone]
    ($(s1.north)+(0.23,0)$)
    .. controls +(0,1.17) and +(0,1.17) ..
    ($(s3.north)+(0.25,0)$);
  \draw[sf skip,draw=strategyarrowtwo]
    ($(s2.north)+(0.02,0)$)
    .. controls +(0,1.72) and +(0,1.72) ..
    ($(s4.north)+(-0.27,0)$);

  % Two visible children plus an ellipsis: the ellipsis is not a child node.
  \foreach \j in {1,...,4}{
    \coordinate (fork\j) at ($(s\j.south)+(0,-0.46)$);
    \node[sf code] (c\j L) at ($(s\j.center)+(-0.43,-1.52)$) {};
    \node[sf code] (c\j R) at ($(s\j.center)+(0.43,-1.52)$) {};
    \node[text=codegreen,font=\small,inner sep=0pt]
      at ($(s\j.center)+(0,-1.52)$) {$\cdots$};
    \draw[draw=codegreen!80,line width=0.85pt]
      (s\j.south) -- (fork\j);
    \draw[sf arrow,draw=codegreen!85,line width=0.85pt]
      (fork\j) -- (c\j L.north);
    \draw[sf arrow,draw=codegreen!85,line width=0.85pt]
      (fork\j) -- (c\j R.north);
  }
\end{tikzpicture}
